% Appendix: Questions and Answers
% User annotations and interactive Q&A records

\appendix
\chapter{问答与注解}

\section{阅读过程中的问题与解答}

% === 用户问答记录 ===
% 本节记录阅读过程中的问答内容

\subsection*{亚组的定义}
\textbf{问}：什么是亚组？

\textbf{答}：亚组（subgroup）是由某个协变量（通常是二元变量 $X_i$）分割出来的子群体。例如按性别、年龄、收入等分割。亚组因果效应衡量的是在该亚组内的平均因果效应。

\subsection*{恒等式 $\tau = \pi_1 \tau_1 + \pi_0 \tau_0$ 的推导}
\textbf{问}：如何推导 $\tau = \pi_1 \tau_1 + \pi_0 \tau_0$？

\textbf{答}：利用指示函数的完备性 $\mathbb{I}(X_i=1) + \mathbb{I}(X_i=0) = 1$，将总体求和拆分为两个亚组求和，再除以对应的亚组样本量即可得到。

\subsection*{为什么在 $H_{0F}$ 下 $\bm{Y}$ 是固定的？}

\textbf{问}：为什么在 Fisher's sharp null ($H_{0F}: Y_i(1) = Y_i(0)$) 下，潜在结果向量 $\bm{Y}$ 是固定的？

\textbf{答}：$H_{0F}$ 断言 treatment 对每个单元都没有效应，即 $Y_i(1) = Y_i(0) = Y_i$ 对所有 $i$。这意味着整个 $\bm{Y}(1) = \bm{Y}(0) = \bm{Y}$ 向量是确定的（固定值），不再有随机性。

在 FRT 中，我们 condition on 这个固定的 $\bm{Y}$，只让治疗分配 $\bm{Z}$ 有随机性。因此，随机化分布可以精确枚举（$M = \binom{n}{n_1}$ 种可能），不依赖任何渐近假设——这就是 FRT"有限样本精确"的本质。

\textit{当你在阅读过程中提出问题时，回答将被记录在此处。}

% === 用户注解 ===
% 使用 \userannotation{问题}{解答} 格式插入

\section{学习笔记}

\textit{在这里记录你的学习心得和思考。}

\endinput
